% CLAUDE REPORT::COSTMEASURE-0001::2026-11-12T21:50Z
% !TEX program = xelatex
\documentclass[11pt]{article}
\usepackage[margin=1in]{geometry}
\usepackage{amsmath,amssymb,amsthm}
\usepackage{xurl}
\newtheorem{theorem}{Theorem}
\newtheorem{corollary}[theorem]{Corollary}
\theoremstyle{remark}\newtheorem*{remark}{Remark}
\newcommand{\Wb}{\overline W}
\newcommand{\status}[1]{\textbf{[#1]}}
\title{CLAUDE COSTMEASURE-0001\\[2pt]\large Cost-measure transfer test for $\Wb_R(N)$, and a 2026 literature check}
\author{CLAUDE (scheduled routine; supplementary, not an LSPV role)}
\date{2026-11-12 13:50 America/Los\_Angeles}
\begin{document}\maketitle
\noindent\textbf{Status.} P versus NP remains \textbf{OPEN}. This entry is additive. It changes no schedule,
allocation, ownership region or root \texttt{main.tex}. It does \emph{not} re-derive items 1--4 of the
stable state in \texttt{research/agents/claude/worst-case-derivation/CONSOLIDATED\_LIVING\_LOG.md}
(bridge collapse, monotonicity, squeeze, portfolio). It adds one red-team test that every proposed
$\Wb_R$ lower bound must pass.

\section*{CORPUS REVIEWED}
The consolidated log and its index of 229 run branches (on branch \texttt{ccr-1695bd63-rv62vk}); the
four routine runs after 17:00 UTC (BRIDGE-0001 at 18:54 and 20:53, WR-BRIDGE-0001 at 19:51); root
\texttt{AGENTS.md}; the LSPV handoffs dated 2026-11-12 under \texttt{research/coordination/}.

\section*{DUPLICATION / CONTRADICTION AUDIT}
All four runs after 17:00 re-derived squeeze, collapse and monotonicity. This repeats stable-state items
2--4. The cause is the one the consolidated log already names: each run starts from \texttt{main},
where the log does not exist. No contradictions were found.

\section{The cost-measure transfer test}
Write $\mathrm{LB}(R,c)$ for a superpolynomial lower bound on $\Wb_R$ under cost measure $c$, and
$\mathrm{PB}$ for its negation (polynomially bounded).
\begin{theorem}\status{PROVED/EXACT}
Let $R'$ simulate $R$ by a map $\sigma$ on refutations. Suppose that for some polynomial $q$,
$c'(\sigma(\pi))\le q(|F|+c(\pi))$. Then $\mathrm{LB}(R',c')\Rightarrow\mathrm{LB}(R,c)$ on the same
CNF families.
\end{theorem}
\begin{proof}
$W_{R'}(F)\le c'(\sigma(\pi^\ast))\le q(|F|+W_R(F))$, where $\pi^\ast$ is an optimal $R$-refutation of $F$.
\end{proof}
\begin{corollary}[Red-team test]\status{PROVED/EXACT}
Suppose a claimed lower bound $\mathrm{LB}(R',c')$ holds on CNF families, for a system $R'$ that
simulates Extended Frege with polynomial cost in $c'$. Then it gives an EF size lower bound, which is a
major open problem. So whenever a lower bound is reported for an ``EF-strong'' system, one of three
things must be true:
\begin{enumerate}
\item the simulation is not cost-polynomial in the measure $c'$ that was bounded;
\item the hard instances are not Boolean CNF encodings (for example, they are non-Boolean or
coefficient-heavy instances);
\item the claimed lower bound is wrong.
\end{enumerate}
\end{corollary}
\begin{remark}[Why this matters for the multisign DAG-hypergraph program]
The minimum-work idea charges $\operatorname{cost}_R$, not proof length. A cost function that is cheap
on extension-like steps (defining new hyperedges or signed nodes) but expensive on some other step can
make $\Wb_R$ superpolynomial while EF-length stays polynomial. Such a lower bound tells us nothing about
algorithms, because a polynomial-time algorithm can pay with length.
Admissibility condition: \emph{$\operatorname{cost}_R$ must be polynomially related, in both
directions, to the bit-size of the derivation in a model-matched encoding.} Only then does
monotonicity (stable item 2) transfer lower bounds downward and bridges upward. The Time--Space
Auditor role should apply this check to every cost model in the project, including
$R_{\mathrm{mix341}}$.
\end{remark}

\section{Literature check (barrier-transfer preflight)}
\begin{itemize}
\item \textbf{ECCC TR26-109}, Fleming--Grosser--Pitassi--Robere, \emph{Provable Reductions in TFNP}
(June 2026). It defines systems $\langle\mathrm{EF},R\rangle$ for TFNP problems $R$. If EF proves $R\in$
FP, then $\langle\mathrm{EF},R\rangle\equiv\mathrm{EF}$. Classification: \emph{useful structural
analogy}. It is a TFNP-indexed version of the provability gap in stable item 6. A SAT solver that is fast
but not EF-provably correct yields a system that may be strictly stronger than EF, so the bridge
$\mathrm{B3}(\mathrm{EF})$ fails exactly when provability fails. Missing transfer theorem: show that
$\mathrm{SAT}\in P$ makes the relevant search problem EF-provably total and correct. This is the
open $H_{\mathrm{PV}}$.
\item \textbf{Theory of Computing 22 (2026)}, Alekseev, \emph{A Lower Bound for Polynomial Calculus
with Extension Rule}. This system simulates EF, so by the Corollary its bound cannot be an EF lower
bound on CNF families. It must fall under case 1 or case 2 above. Our reading (not re-verified from
the paper): the hard instances are non-Boolean or coefficient-heavy. Classification: \emph{failed
transfer for the bridge}, but it is a good worked example of the cost-measure test.
\status{CANDIDATE reading}
\item \textbf{ECCC TR26-055}, Davis--Robere: unchanged; already logged as an analogy.
\end{itemize}
No 2026 source gives an unconditional EF lower bound or an optimal proof system. Stable item 6 stays
\status{OPEN}.

\section*{ADOPTED / REJECTED / RETIRED}
Adopted: the admissibility condition on $\operatorname{cost}_R$. Rejected: using lower bounds under
non-size-related cost measures as evidence for $\Wb_R$ hardness. Nothing was retired.

\section*{STALE OR OVERSTATED CLAIMS FOUND}
None new. Process defect: the consolidated log is still not on \texttt{main}, so runs keep repeating work.

\section*{WHAT CHANGED RELATIVE TO ROOT MAIN}
Nothing. This routine owns no region of \texttt{main.tex}.
\end{document}
% END_OF_VERIFIED_REPORT::CLAUDE::COSTMEASURE-0001
